%! Introducing Vectors — Unit 1, Lesson 1.0 — Warm-Up (Answer Key)
\documentclass[10pt]{article}
\usepackage{linalg-article}
\usepackage{linalg-key}   % pulls in -boxes; do NOT also load -boxes
\newcommand{\TallMath}[1]{$\displaystyle #1\rule[-1.4em]{0pt}{3.2em}$}

\begin{document}

\pageheader{Unit 1, Lesson 1.0}{Warm-Up --- Answer Key}
\namedateperiod

\vspace{0.12in}

\begin{notesbox}{Getting Ready: Skills We'll Use Today}
\small These three skills power everything in today's lesson. Answer quickly.

\vspace{4pt}
\textbf{1. Signed-number arithmetic.} Evaluate.
\begin{enumerate}[label=(\alph*), leftmargin=2.0em, itemsep=6pt]
  \item $3 + (-7) = {}$\ans{$-4$}
  \item $-4 - (-9) = {}$\ans{$5$}
  \item $2(-5) = {}$\ans{$-10$}
\end{enumerate}

\vspace{6pt}
\textbf{2. The Pythagorean theorem.} A right triangle has legs of length $6$ and $8$.
How long is the hypotenuse? Show the step.
\ansline{$\sqrt{6^2+8^2}=\sqrt{36+64}=\sqrt{100}=10$}

\vspace{4pt}
\textbf{3. Reading the plane.} On the grid, point $P$ is marked.
\begin{center}
\begin{tikzpicture}[scale=0.5]
  \draw[step=1, linegray!60, thin] (-0.4,-0.4) grid (5.4,4.4);
  \draw[->, thick, charcoal] (0,0)--(5.6,0) node[right]{\scriptsize $x$};
  \draw[->, thick, charcoal] (0,0)--(0,4.6) node[above]{\scriptsize $y$};
  \foreach \x in {1,2,3,4,5} \node[below, font=\scriptsize] at (\x,-0.1){\x};
  \foreach \y in {1,2,3,4} \node[left, font=\scriptsize] at (-0.1,\y){\y};
  \fill[burgundy] (5,3) circle (3pt) node[above right]{\small $P$};
\end{tikzpicture}
\end{center}
What are the coordinates of $P$? \; $\big(\,\ans{$5$},\ \ans{$3$}\,\big)$
\end{notesbox}

\begin{teachernote}
Item 2 previews the magnitude formula exactly: today's $\|\mathbf{v}\|=\sqrt{v_1^2+v_2^2}$
\emph{is} the Pythagorean theorem. Item 3 rehearses reading an ordered pair, which becomes
reading a vector's components. If item 1c is shaky, revisit scalar multiplication slowly.
\end{teachernote}

\end{document}
